\section{Abschlussphase} 
\label{sec:Abschlussphase}
Zum Ende der Dokumentation wird das Projekt in der Retroperspektive betrachtet. In einem Soll-Ist-Vergleich wird geprüft, ob alle
Anforderungen beachtet und die Zeitplanung eingehalten wurde. Abschließend wird das Projekt mit einem
persönlichen Fazit bewertet.


\subsection{Soll-Ist-Vergleich}
\label{sub:Soll-Ist-Vergleich}
Alle Anforderungen an die Erweiterung konnten erfüllt werden. Die Art der Aktualisierung der Status-Anzeigen ist, nach meinem
Empfinden, unzufrieden gelöst. Siehe \VerweisKL{sub:Datenübertragung}. Dies musste jedoch zur Einhaltung des zeitlichen Rahmens
vorerst hingenommen werden.\\
\endgroup
In \TabelleK{tab:Zeitplanung} werden der angenommene und der tatsächliche benötigte Zeitaufwand verglichen. Wie man erkennen
kann, dauerte die Durchführung des Proof-of-Concept wesentlich länger als geplant. Dies lag daran, dass das Framework
\Fachbegriff{Atmosphere} nicht in einem Umfang wie gewünscht dokumentiert ist. Auch da es stetig Änderungen am Framework gibt,
sind Demo-Programmcodes häufig veraltet gewesen. Der Mehraufwand reduzierte jedoch die benötigte Zeit für die Einarbeitung ins
Framework und das Erstellen eines Grundgerüstes. Der restliche Mehraufwand verkleinerte den Puffer für Anforderungsänderungen. Es
ist zu sehen, dass dieser gänzlich ausgeschöpft wurde. Zunächst für seinen geplanten Sinn, in der markierten Zeile, für die
Umsetzung einer Anforderungsänderung. \\
Der übrige Puffer konnte dann für eine bessere \Abkuerzung{UI}-Umsetzung und eine umfangreichere Qualitätssicherung genutzt
werden.

\newTabelle{Ursprüungliche Zeitplanung und tatsächlich benötigte Zeit im Vergleich}{tab:Zeitplanung}{Zeitplanung.tex}
Erklärung der unterstrichenen Zeile: Die aus dem Projektantrag stammende Bezeichnung "`Datenbank und Entität anpassen"' ist
aufgrund einer Anforderungsänderung nicht mehr zutreffend und wurde angepasst.

\subsection{Persönliches Fazit}
\label{sub:Fazit}
Das Projektthema, die Realisierung eines Benachrichtigungssystems, hat mich seit Nennung im \Abkuerzung{VMS}-Projekt als
Anforderung gereizt. Es ist nicht nur ein schönes IHK-Abschlussprojekt, es ist auch ein sehr eigenständiges und abgrenzbares
Feature. Ich konnte mich dabei fachlich weiterbilden, da ich mich in meinem Projektthema nicht auf bereits Beherrschtes beschränkt
habe. Darüber hinaus war eine persönliche Entwicklung zum besseren Anwendungsentwickler zu sehen. Ich freue mich auch künftig am
Projekt \Abkuerzung{VMS} mitwirken zu können, sowie das Thema der Erweiterung weiter zu begleiten.
